% =========================================================
% Proof: Every Infinite SB Path Converges to Its Target
% =========================================================

\newpage
\phantomsection
\label{prf:sb-path-limit}

\begin{remark*}[Return]
\hyperref[cor:sb-path-limit]{Return to Corollary~\ref*{cor:sb-path-limit}.}
\end{remark*}

\begin{theorem*}[Every infinite Stern--Brocot path converges to its target]
Let $x > 0$ and let $(I_n)$ be the SB refinement sequence targeting $x$.
Then $a_n/b_n \to x$ and $c_n/d_n \to x$.
\end{theorem*}

\begin{proof}[Proof (Professional Standard)]
By \hyperref[prop:sb-unique-point]{Proposition~\ref*{prop:sb-unique-point}},
there is a unique $\xi \in \bigcap I_n$ and both endpoint sequences
converge to $\xi$.  By the construction rule, $x \in I_n$ for every $n$
(at each step, the endpoint on the same side as the mediant relative to
$x$ is replaced, keeping $x$ in the new interval).  Therefore
$x \in \bigcap I_n$.  By uniqueness, $\xi = x$.  Hence both endpoint
sequences converge to $x$.
\end{proof}

\begin{proof}[Proof (Detailed Learning)]
\textbf{Step 1: Apply the Unique Point Proposition.}
By \hyperref[prop:sb-unique-point]{Proposition~\ref*{prop:sb-unique-point}},
$\bigcap I_n = \{\xi\}$ for some $\xi \in \mathbb{R}$, and both
$(a_n/b_n)$ and $(c_n/d_n)$ converge to $\xi$.

\textbf{Step 2: Verify $x \in I_n$ for every $n$.}
The construction rule is: at each step, compare $x$ to the mediant $m$.
If $x < m$, replace the right endpoint by $m$ (so the new interval is
$[a_n/b_n, m]$ and $x < m$ means $x$ is still in the left part).
If $x > m$, replace the left endpoint (similarly).
If $x = m$, the process terminates (but we are in the infinite-path
case, so $x \neq m$ at every step).
In either case, $x$ lies in the updated interval.  By induction, $x \in I_n$
for all $n$.

\textbf{Step 3: Conclude $\xi = x$.}
Since $x \in I_n$ for all $n$, we have $x \in \bigcap I_n = \{\xi\}$.
Therefore $x = \xi$.

\textbf{Step 4: State the convergence.}
Both sequences converge to $\xi = x$:
\[
\lim_{n\to\infty} \frac{a_n}{b_n} = x,
\qquad
\lim_{n\to\infty} \frac{c_n}{d_n} = x.
\qedhere
\]
\end{proof}

\begin{remark*}[Proof structure]
A one-step corollary: the unique point in $\bigcap I_n$ is $\xi$; the
target $x$ lies in every $I_n$ by the construction rule; uniqueness forces
$\xi = x$.  All the work is done by the Unique Point Proposition.
\end{remark*}

\begin{remark*}[Dependencies]
\hyperref[prop:sb-unique-point]{Nested SB intervals contain a unique point},
the Stern--Brocot construction rule.
\end{remark*}
